%% kapa: exam=2 round=31 year=2020 period=3 subject=compensation_law subject_ko=감정평가 및 보상법규 seq=3 number=3 points=20
%% source: Q-Net 감정평가사 기출문제 (cst00309) artlSeq=5207673 file=raw/1008/5207673/3교시_감정평가 및 보상법규.pdf
%% conversion: kordoc 4.24.51 -> tools/gen2.py (auto); review=proofread-gpt-astra; flags=none
\begin{longquestion}{【문제 3】}
甲과 乙은 감정평가사 자격이 없는 공인회계사로서, 甲은 A주식회사의 부사장 겸 본부장이고 乙은 A주식회사의 상무의 직에 있는 자이다. 甲과 乙은 A주식회사 대표 B로부터 서울 소재의 A주식회사 소유 빌딩의 부지를 비롯한 지방에 있는 같은 회사 전 사업장 물류센터 등 부지에 대한 자산재평가를 의뢰받고, 회사의 회계처리를 목적으로 부지에 대한 감정평가 등 자산재평가를 실시하여 그 결과 평가대상 토지(기존의 장부상 가액 3천억원)의 경제적 가치를 7천억원의 가액으로 표시하고, 그 대가로 1억 5,400만원을 받았다. 이러한 甲과 乙의 행위는「감정평가 및 감정평가사에 관한 법률」상의 감정평가업자의 업무에 해당하는지 여부에 관하여 논하시오. (20점) \par
\end{longquestion}
